\documentclass[10pt,twocolumn]{article}

% ---------- PACKAGES ----------
\usepackage[a4paper,margin=1.8cm,columnsep=0.6cm]{geometry}
\usepackage{times}
\usepackage[T1]{fontenc}
\usepackage[utf8]{inputenc}
\usepackage{amsmath,amssymb}
\usepackage{graphicx}
\usepackage{caption}
\usepackage{subcaption}
\usepackage{booktabs}
\usepackage{listings}
\usepackage{xcolor}
\usepackage{hyperref}
\usepackage{titlesec}
\usepackage{abstract}
\usepackage{enumitem}
\usepackage{cite}

% ---------- LISTINGS STYLE ----------
\lstdefinestyle{code}{
    basicstyle=\ttfamily\scriptsize,
    keywordstyle=\color{blue},
    commentstyle=\color{gray},
    stringstyle=\color{teal},
    numbers=left,
    numberstyle=\tiny\color{gray},
    breaklines=true,
    frame=single,
    showstringspaces=false,
    tabsize=2
}
\lstset{style=code}

% ---------- TITLE FORMAT ----------
\titlespacing*{\section}{0pt}{8pt}{4pt}
\titlespacing*{\subsection}{0pt}{6pt}{2pt}

\setlength{\abstitleskip}{-10pt}
\renewcommand{\abstractnamefont}{\normalfont\bfseries}
\renewcommand{\abstracttextfont}{\normalfont\itshape}

\title{\Large \textbf{Parallel Simulation of the Damped Wave Equation:\\
A Comparative Study Using OpenMP, MPI, and CUDA}}

\author{
\begin{tabular}[t]{ccc}
Allione Paolo & Frulla Edoardo & Porro Michele \\
\small 345215 & \small 343236 & \small 342079
\end{tabular}\\[4pt]
\small Quantum Engineering -- High Performance Computing Course\\
\small Academic Year 2025--26
}



\date{}

\begin{document}

\twocolumn[
\begin{@twocolumnfalse}
\maketitle
\begin{abstract}
This report presents the design, implementation, and performance evaluation of a numerical
simulator for the two-dimensional damped wave equation, developed progressively using three
different parallel programming paradigms: OpenMP, MPI, and CUDA. The first part of the
project implements a shared-memory parallel solver based on a finite-difference discretization
of the governing partial differential equation, producing a sequence of grayscale frames later
assembled into a video. The second part extends the simulator to execute multiple independent
wave simulations concurrently using a hybrid MPI+OpenMP approach, modeling single-source and
multi-source (interfering) wave phenomena. The third part introduces GPU-accelerated
post-processing with CUDA, where each simulation frame is recolored to emphasize wave maxima,
minima, and propagation fronts. For each part, the numerical method, the parallelization
strategy, and the experimental results -- including correctness validation and performance
scaling -- are discussed.
\end{abstract}
\vspace{0.3cm}
\end{@twocolumnfalse}
]

% =====================================================================
\section{Introduction}
\label{sec:intro}
% =====================================================================
The damped wave equation is a second-order partial differential equation that models the
propagation of a perturbation through a medium subject to energy dissipation, such as
friction, viscosity, or internal resistance. In two spatial dimensions, the equation is
written as

\begin{equation}
\frac{\partial^2 u}{\partial t^2} + \gamma \frac{\partial u}{\partial t} = c^2 \nabla^2 u,
\label{eq:wave}
\end{equation}

where $u(x,y,t)$ represents the wave height on the plane $(x,y)$ at time $t$, $\gamma$ is the
damping coefficient, and $c$ is the propagation speed of the wave. Compared to the ideal,
non-dissipative wave equation, the term $\gamma \, \partial u / \partial t$ progressively
attenuates the amplitude of the wave as it propagates.

Solving Equation~\eqref{eq:wave} numerically requires discretizing both the spatial domain
(grid spacing $\Delta x = \Delta y$) and the temporal domain (time step $\Delta t$), replacing
the continuous function $u(x,y,t)$ with a discrete field $u_{i,j}^{n}$, where $(i,j)$ are the
grid indices and $n$ is the discrete time step. Using a second-order central difference
approximation for both time derivatives and a five-point stencil for the Laplacian operator,
the update rule for the wave field at time step $n+1$ is obtained as

\begin{equation}
\nabla^2 u_{i,j}^{n} \approx
\frac{u_{i+1,j}^{n} + u_{i-1,j}^{n} + u_{i,j+1}^{n} + u_{i,j-1}^{n} - 4u_{i,j}^{n}}{\Delta x^2},
\label{eq:laplacian}
\end{equation}

\begin{equation}
u_{i,j}^{n+1} =
\frac{2u_{i,j}^{n} - \left(1 - \dfrac{\gamma \Delta t}{2}\right)u_{i,j}^{n-1}
+ c^2 \Delta t^2 \nabla^2 u_{i,j}^{n}}
{1 + \dfrac{\gamma \Delta t}{2}}.
\label{eq:update}
\end{equation}

Equation~\eqref{eq:update} forms the computational core of all three parts of this project,
and it is evaluated, frame by frame, for every grid point of an $M \times M$ matrix, over $N$
simulation steps. This report is organized into three sections, each corresponding to one part
of the assignment: Section~\ref{sec:openmp} describes the shared-memory OpenMP implementation
of the base simulator; Section~\ref{sec:mpi} presents the MPI-based extension that executes
three concurrent simulations, including single- and multi-source wave interference; and
Section~\ref{sec:cuda} discusses the CUDA-accelerated frame colorization stage. Each section
introduces the corresponding sub-problem and is followed by a placeholder for the experimental
results obtained on the target HPC cluster.

% =====================================================================
\section{Part 1 -- Shared-Memory Simulation with OpenMP}
\label{sec:openmp}
% =====================================================================

\subsection{Problem Description}
The first part of the project requires the implementation of a sequential-to-parallel
simulator of the damped wave equation introduced in Section~\ref{sec:intro} (Eq.~\ref{eq:wave}),
parallelized using OpenMP. Given the physical parameters $\gamma$ (damping) and $c$
(propagation speed), and the numerical parameters $\Delta t$, $\Delta x$, the grid size $M$,
and the number of time steps $N$, the simulator must compute the evolution of the wave height
matrix $u_{i,j}^{n}$ starting from an initial condition where the matrix is uniformly zero
except for an initial impulse applied at a boundary point $(i_0, j_0)$ at $n=0$.

\subsection{Numerical Method and Implementation}
At each time step, the simulator applies the explicit update rule of
Equation~\eqref{eq:update}, which requires keeping in memory the wave fields at the two
previous time steps ($u^{n}$ and $u^{n-1}$) to compute $u^{n+1}$. The computation is organized
as a loop over the $N$ requested frames; for each frame, the update is applied independently
to every interior grid point $(i,j)$, which makes the spatial loop embarrassingly parallel and
well suited to a \texttt{\#pragma omp parallel for} directive, typically applied with
\texttt{collapse(2)} on the nested $i,j$ loops to maximize the available parallelism. Boundary
conditions are imposed by clamping or reflecting the edges of the matrix.

After each time step is computed, the resulting matrix is normalized to the $[0,255]$ range
and written to disk as a grayscale image in PGM (Portable Graymap) format, following the P5
binary encoding, into a dedicated \texttt{sim} folder, with filenames of the form
\texttt{frame\_\%05d.pgm}. Once all $N$ frames have been generated, the FFmpeg tool is used to
assemble the image sequence into an MP4 video, with the frame rate chosen consistently with the
simulation time step $\Delta t$ (i.e., $\text{fps} = 1/\Delta t$, suitably scaled for a
perceptually smooth playback).

\subsection{Parallelization Strategy}
% TODO (optional): briefly justify your scheduling choice (static/dynamic),
% number of threads tested, and any synchronization/false-sharing considerations.

\subsection{Results}
\textit{[Insert here: simulation parameters used (Table~\ref{tab:part1-params}), example
frames showing wave propagation at different time steps (Figure~\ref{fig:part1-frames}),
execution time vs.\ number of OpenMP threads, and the resulting speedup/efficiency curves.]}

\begin{table}[h]
\centering
\caption{Simulation parameters -- Part 1.}
\label{tab:part1-params}
\begin{tabular}{ll}
\toprule
Parameter & Value \\
\midrule
$\gamma$ (damping)        & \textit{TBD} \\
$c$ (propagation speed)   & \textit{TBD} \\
$\Delta t$                & \textit{TBD} \\
$\Delta x$                & \textit{TBD} \\
$M$ (grid size)           & \textit{TBD} \\
$N$ (number of frames)    & \textit{TBD} \\
\bottomrule
\end{tabular}
\end{table}

\begin{figure}[h]
\centering
% \includegraphics[width=\linewidth]{frames_openmp.png}
\caption{Example frames of the wave simulation at different time steps. \textit{[Placeholder -- insert figure]}}
\label{fig:part1-frames}
\end{figure}

\begin{figure}[h]
\centering
% \includegraphics[width=\linewidth]{speedup_openmp.png}
\caption{Speedup and efficiency as a function of the number of OpenMP threads. \textit{[Placeholder -- insert figure]}}
\label{fig:part1-speedup}
\end{figure}

% =====================================================================
\section{Part 2 -- Multi-Simulation Extension with MPI}
\label{sec:mpi}
% =====================================================================

\subsection{Problem Description}
The second part of the project extends the OpenMP simulator of Section~\ref{sec:openmp} to
the distributed-memory setting using MPI, with the goal of executing \emph{three independent
simulations simultaneously} on an HPC cluster. Each simulation continues to exploit
intra-node, shared-memory parallelism through OpenMP, resulting in a hybrid MPI+OpenMP
implementation. The three simulations differ in their initial conditions:

\begin{enumerate}[label=(\roman*)]
    \item \textbf{Simulation 1}: identical to the base case of Part~1, with a single impulse
    source at $(i_0,j_0)$ and $t=0$.
    \item \textbf{Simulation 2}: two independent wave sources, located at two distinct points
    of the plane, sharing the same damping $\gamma$ and propagation speed $c$, both activated
    at $t=0$. This configuration is designed to highlight wave interference patterns arising
    from the superposition of the two wavefronts.
    \item \textbf{Simulation 3}: identical to Simulation 2, but the second source is activated
    at a later time $t = t_{\text{start}} > 0$, producing an asymmetric interference pattern.
\end{enumerate}

\subsection{Parallelization Strategy}
The three simulations are mapped onto separate MPI processes (or groups of processes),
allowing them to run concurrently and independently, since there is no data dependency between
them. Each MPI process executes its own instance of the OpenMP-parallelized time-stepping loop
described in Section~\ref{sec:openmp}, writing its output frames to a dedicated folder
(\texttt{sim1}, \texttt{sim2}, \texttt{sim3} respectively). This two-level parallelization
scheme -- MPI across simulations and OpenMP within each simulation -- allows efficient use of
a multi-node cluster where each node hosts one MPI process exploiting all of its available
cores via OpenMP threads.

% TODO (optional): describe communication pattern (if any communication between
% processes is needed, e.g., for synchronized output or load balancing), and how
% the two interfering sources are summed/composed onto the shared grid in
% Simulations 2 and 3.

\subsection{Results}
\textit{[Insert here: simulation parameters for the three runs (Table~\ref{tab:part2-params}),
example frames illustrating single-source propagation and the interference patterns of
Simulations 2 and 3 (Figure~\ref{fig:part2-frames}), and the measured execution time/scalability
of the hybrid MPI+OpenMP implementation as a function of the number of MPI processes and OpenMP
threads.]}

\begin{table}[h]
\centering
\caption{Simulation parameters -- Part 2.}
\label{tab:part2-params}
\begin{tabular}{lccc}
\toprule
Parameter & Sim. 1 & Sim. 2 & Sim. 3 \\
\midrule
$\gamma$               & \textit{TBD} & \textit{TBD} & \textit{TBD} \\
$c$                    & \textit{TBD} & \textit{TBD} & \textit{TBD} \\
Source(s) position     & \textit{TBD} & \textit{TBD} & \textit{TBD} \\
$t_{\text{start}}$ (2nd source) & -- & 0 & \textit{TBD} \\
\bottomrule
\end{tabular}
\end{table}

\begin{figure}[h]
\centering
% \includegraphics[width=\linewidth]{frames_mpi.png}
\caption{Comparison of single-source propagation and multi-source interference patterns. \textit{[Placeholder -- insert figure]}}
\label{fig:part2-frames}
\end{figure}

\begin{figure}[h]
\centering
% \includegraphics[width=\linewidth]{scaling_mpi.png}
\caption{Execution time and scalability of the hybrid MPI+OpenMP implementation. \textit{[Placeholder -- insert figure]}}
\label{fig:part2-scaling}
\end{figure}

% =====================================================================
\section{Part 3 -- GPU-Accelerated Frame Colorization with CUDA}
\label{sec:cuda}
% =====================================================================

\subsection{Problem Description}
The third part of the project focuses on improving the visual quality of the simulation output
by transforming each grayscale frame into a colored RGB frame, with the goal of clearly
highlighting the maxima and minima of the wave field while smoothly fading the intermediate
(``calm'') regions, which are colored with a neutral, bright color (e.g., white) to make the
propagating wavefronts stand out. The colorization step is implemented on the GPU using CUDA,
since the mapping from a scalar grayscale value to an RGB triplet is an embarrassingly
parallel, per-pixel operation, making it an ideal candidate for massive data-parallel execution.

Each colored frame is saved in PPM (Portable Pixmap) format, using the ASCII P3 encoding, in
which every pixel is represented as a triplet of values $(R,G,B) \in [0,255]^3$. The original
grayscale matrix, computed by the OpenMP/MPI simulation core of Parts~1--2, is left unmodified
and is used to continue the time-stepping of the simulation; only a colored copy of the frame
is produced for visualization purposes.

\subsection{Color Mapping and CUDA Kernel}
A custom transfer function $f: [0,255] \rightarrow [0,255]^3$ is defined to map a grayscale
intensity (representing the wave amplitude, after normalization) to an RGB color. The mapping
is designed so that:
\begin{itemize}
    \item values near the neutral/zero level of the wave are mapped to a light, low-saturation
    color (white or near-white), so that calm regions of the plane remain visually unobtrusive;
    \item values approaching the maximum amplitude are mapped to a bright, highly saturated
    color (e.g., red), and values approaching the minimum amplitude to a different bright color
    (e.g., blue), so that crests and troughs of the wave are immediately distinguishable;
    \item intermediate values (the rising and falling fronts of the wave) are mapped through a
    smooth color gradient interpolating between the neutral color and the extremal colors, in
    order to avoid visual banding artifacts.
\end{itemize}

This transfer function is implemented as a CUDA kernel that is launched with one thread per
pixel of the $M \times M$ grayscale matrix; each thread independently reads the corresponding
grayscale value, applies the color transfer function, and writes the resulting RGB triplet to
the output buffer, which is then transferred back to host memory and serialized to the PPM
file. Because each pixel is processed independently, the kernel requires no inter-thread
synchronization or shared memory beyond what is used for performance optimization (e.g.,
coalesced global memory access, or optional use of constant/texture memory to store
lookup tables for the color gradient).

% TODO (optional): specify CUDA grid/block configuration, memory transfer
% strategy (e.g., overlapping computation/transfer with streams), and whether
% a precomputed color lookup table (LUT) was used to accelerate the mapping.

\subsection{Results}
\textit{[Insert here: example colorized frames showing the maxima/minima highlighting and
front fading effect (Figure~\ref{fig:part3-frames}), comparison between CPU and GPU
colorization time per frame, and overall GPU kernel performance/throughput analysis
(Figure~\ref{fig:part3-perf}).]}

\begin{figure}[h]
\centering
% \includegraphics[width=\linewidth]{frames_cuda.png}
\caption{Colorized simulation frames highlighting wave maxima, minima, and fading fronts. \textit{[Placeholder -- insert figure]}}
\label{fig:part3-frames}
\end{figure}

\begin{figure}[h]
\centering
% \includegraphics[width=\linewidth]{perf_cuda.png}
\caption{CUDA kernel execution time vs.\ grid size, and comparison against a CPU-only colorization baseline. \textit{[Placeholder -- insert figure]}}
\label{fig:part3-perf}
\end{figure}

% =====================================================================
\section{Conclusions}
% =====================================================================
This report described the incremental development of a damped wave equation simulator across
three parallel computing paradigms. The OpenMP implementation established a correct and
efficient shared-memory baseline; the MPI extension enabled the concurrent execution of
multiple independent simulations, including the study of wave interference phenomena, in a
hybrid MPI+OpenMP fashion suitable for multi-node HPC clusters; and the CUDA stage demonstrated
how GPU acceleration can be effectively applied to a massively data-parallel post-processing
task, improving the visual interpretability of the simulation results without affecting the
underlying numerical computation. \textit{[Insert here a brief summary of the quantitative
findings: best speedups obtained, scalability limits observed, and overall lessons learned.]}

% =====================================================================
% \section*{References}
% Add bibliography here if required by the course, e.g. via \texttt{thebibliography}
% or a \texttt{.bib} file with \texttt{natbib}/\texttt{biblatex}.
% =====================================================================

\end{document}
